\section{Lambek calculus in linguistics}
In linguistics, the Lambek calculus is best understood as part of the broader tradition of
categorial grammar. Categories behave like types, and grammatical composition is governed by
logical operations rather than by phrase-structure rules. This makes the calculus attractive for
analysing phenomena where lexical information and local combinatorics are central.

Type-logical grammar developed this perspective in depth and connected it to proof theory,
formal semantics, and computational linguistics \citep{Moortgat1997,Morrill1994}. Closely related
systems, especially Combinatory Categorial Grammar (CCG), have been particularly successful in
wide-coverage parsing and semantic interpretation \citep{Steedman2000,ClarkCurran2007}.
Although CCG is not identical to the Lambek calculus, the family resemblance is close enough
that probabilistic work on CCG provides a useful source of intuition for the present project.

